\subsection*{Cancellation Laws}
\label{subsec:cancellation-laws}

\begin{proposition}[Cancellation Laws in a Group]
\label{prop:group-cancellation}
Let $G$ be a group and let $a,b,c \in G$.
If $ab = ac$, then $b=c$.
If $ba = ca$, then $b=c$.
\hyperref[prf:group-cancellation]{Go to proof.}
\end{proposition}

\begin{remark*}[Standard quantified statement]
\[
\operatorname{Group}(G,\star)
\Longrightarrow
\forall a,b,c\in G,\;
(a\star b=a\star c \Rightarrow b=c)
\land
(b\star a=c\star a \Rightarrow b=c).
\]
\end{remark*}

\begin{remark*}[Theorem predicate reading]
In a group, a common left or right factor can be cancelled from an equation.
\end{remark*}

\begin{remark*}[Negated quantified statement]
\[
\operatorname{Group}(G,\star)
\land
\exists a,b,c\in G,\;
(a\star b=a\star c \land b\neq c)
\lor
(b\star a=c\star a \land b\neq c)
\Longrightarrow \bot.
\]
\end{remark*}

\begin{remark*}[Interpretation]
Cancellation is the algebraic effect of having inverses.
\end{remark*}

\begin{dependencies}
\begin{itemize}
  \item \hyperref[def:algebraic-group]{Group}
  \item \hyperref[def:inverse-element]{Inverse Element}
\end{itemize}
\end{dependencies}

\begin{remark*}[Proof TODO]
Regenerate this proof as a one-proof-per-file artifact.
\end{remark*}
